\documentclass[11pt]{article}
\usepackage[a4paper,margin=1in]{geometry}
\usepackage{amsmath,amssymb,amsthm}
\usepackage{booktabs}
\usepackage{graphicx}
\usepackage{xcolor}
\usepackage{microtype}
\usepackage{hyperref}
\hypersetup{colorlinks=true,linkcolor=blue,urlcolor=blue,citecolor=blue}
\newcommand{\rt}[1]{\texttt{#1}}
\newcommand{\fn}[1]{\texttt{#1}}
%% This report is mostly filenames, which are long unbreakable words; stretch
%% the interword space rather than either running into the margin or hyphenating
%% a path at an arbitrary character.
\sloppy
\newcommand{\wcc}{\mathrm{wcc}}
\newcommand{\ifgraphic}[2]{\IfFileExists{#1}{\includegraphics[width=#2]{#1}}%
  {\fbox{\parbox{#2}{\centering\small figure \texttt{\detokenize{#1}} pending}}}}
\newcommand{\iftable}[1]{\IfFileExists{#1}{\input{#1}}%
  {\fbox{\parbox{0.9\textwidth}{\centering\small table
  \texttt{\detokenize{#1}} pending}}}}

\title{\textbf{Report R19 --- the paper figures: the unitary census, $D_{\max}$,
and how big the sectors actually are}\\
\large Krylov sector analysis of quantum cellular automata}
\author{Tier 1 analysis}
\date{2026-08-13 (rev.~2; rev.~1 2026-08-12)}

\begin{document}
\maketitle

\begin{abstract}
This report accompanies \fn{notebooks/paper\_figures.ipynb}, an editable rebuild
of R2's Figure~1 --- the number of Krylov sectors against $N$ at \rt{obc0} ---
together with the matching $D_{\max}$ figure and the distribution of sector sizes
for the two exponentially fragmented rules W156 and W108. Revision~2 widens the
scaling figures from five tracks to \textbf{all sixteen unitary rules}, ergodic
ones included, which turns them from an illustration into a census: every rule
built from $I$ and $V$ alone does exactly one of four things, and the two trivial
behaviours --- one sector of $2^N$, or $2^N$ sectors of one --- are the
\emph{frame} of both panels rather than omissions from them. They are also the
only rules that saturate $n_\wcc D_{\max} \ge 2^N$. The census is also
\emph{uniform}: the ergodic rules had been left at $N \le 16$ by an early-exit in
the sweep, and \fn{extend\_unitary} closed that gap with the flip reduction, so
all sixteen rules now run to $N = 22$ on the same grid --- $93$ new units, each
gated on the new code first reproducing every unit both sweeps had already
stored for that rule. Four things are worth recording beyond the pictures.
First, every curve is an \textbf{exact integer law}, checked at every computed
$N$ rather than fitted --- $44$ of the $48$ (rule, series) slots, the other four
being genuinely non-closed-form --- so the annotations the figures carry are
certificates (Table~\ref{tab:r19laws}), and all of them survived the extension
to $N=22$.
Second, \textbf{the vacuum alone decides whether some rules fragment}: reflection
is an \rt{obc0} symmetry and collapses sixteen rules to twelve curves, but the
spin flip is not, and it puts the polynomially fragmented W60 and W102 opposite
the completely ergodic W195 and W153. Third, W105 and W150 --- routinely quoted
as one entry ``$150/105$'' --- differ at $N\equiv1\pmod 4$ and at no other $N$,
in \emph{both} series at once, which is also the residue class at which W105 has
no frozen state at all. Fourth, the size distribution separates the two
exponential rules in a way neither scaling figure does: W156's is
\textbf{unimodal}, W108's \textbf{monotone decreasing} with $33\,552$ singletons
out of $170\,625$ sectors at $N=21$, and in both the sector a random
computational-basis state lands in is far smaller than $D_{\max}$.
\end{abstract}

\tableofcontents

\section{What this report is}

R2 introduced the two scaling figures and R9 the growth-base plane built on
them. What was missing was a version of the figures that can be \emph{edited}:
the package figure code in \fn{scaling/figure.py} draws all 256 rules on a fixed
template, which is the right tool for a sweep and the wrong one for a paper
plate. \fn{notebooks/paper\_figures.ipynb} is the replacement. All the drawing
code is inline in the notebook, all the adjustable parameters are collected in
one cell at the top, and no cell calls the package's figure module.

What the notebook does \emph{not} own is the arithmetic. The closed forms
annotated on the curves live in \fn{scaling/paper\_figures.py}, and the notebook
imports and re-verifies them before it draws anything. That split is deliberate:
a law is a fact about the data, an axis limit is a preference, and only the
second should be casually editable.

\paragraph{Outputs.} \fn{figures/paper/}: \fn{fig\_paper\_sectors\_obc0.pdf},
\fn{fig\_paper\_dmax\_obc0.pdf}, \fn{fig\_paper\_scaling\_obc0.pdf} (the two as
panels (a) and (b) of one plate), \fn{fig\_paper\_sizehist\_obc0.pdf} and, new in
this revision, \fn{fig\_paper\_frozen\_obc0.pdf}. Each with a 300\,dpi PNG beside
it.

\section{The census: sixteen rules, four behaviours}
\label{sec:census}

A rule is unitary exactly when none of its four local channels is a reset, i.e.\
when its tuple lies in $\{I,V\}^4$; there are sixteen of them. Revision~1 of this
report drew five of them. The other eleven were not missing because they are
uninteresting but because of a loader: \fn{summary.load\_series} drops units
flagged ergodic, which is correct when fitting a growth base and returns an
\emph{empty} series for seven rules. The notebook now uses its own loader and
they are all present.

\iftable{tab_r19_census_obc0.tex}

\paragraph{The four classes.} Every unitary rule falls into exactly one, with no
intermediate case anywhere in the computed range:
\begin{itemize}
  \item \textbf{ergodic} --- W51, W57, W99, W147, W153, W195 give literally one
    sector containing all $2^N$ states, at every computed $N$; the six series are
    not merely similar but identical. W54 is the same with one exception: the
    all-zero state is frozen, because $r_{00} = I$, so it has $2$ sectors of
    sizes $2^N-1$ and $1$.
  \item \textbf{polynomial} --- W60, W102, W105, W150 fragment into $\Theta(N)$
    sectors that are individually enormous. W60's largest is exactly $2^{N-1}$,
    half the Hilbert space.
  \item \textbf{exponential} --- W108, W156, W198, W201 fragment into $a^N$
    sectors with $1 < a < 2$.
  \item \textbf{frozen} --- W204 is the identity circuit: $2^N$ sectors of one
    state each.
\end{itemize}

\paragraph{The trivial classes are the frame, not the margin.} The hyperbola
$n_\wcc D_{\max} \ge 2^N$ (R16) holds with equality exactly when every sector has
the same size. Two behaviours manage that --- one sector of $2^N$ and $2^N$
sectors of one --- and they are precisely the two that do not fragment. The last
column of Table~\ref{tab:r19census} is that ratio: exactly $1$ for the six
ergodic rules and for the identity, $2$ for W54, and up to $1.3\times10^3$ for
W201. So the ergodic and identity curves in Figure~\ref{fig:r19scaling} are not
extra data points on the same footing as the rest; they are the boundary of the
region, and fragmentation is what happens strictly inside it.

\section{Two symmetries, one of which the boundary breaks}
\label{sec:symmetry}

Sixteen rules do not give sixteen curves, and the reason is worth stating
precisely because the two candidate symmetries behave oppositely at \rt{obc0}.

\paragraph{No symmetry is a source of data.} Reflection \emph{is} a symmetry of
these series, and the sweep driver offers a representatives-only mode, so
``W198 $=$ W156'' could in principle have been a substitution rather than a
measurement. It is not, and the report does not ask the reader to take that on
trust. Every unitary rule carries its own records with its own runtimes ---
W156 and W198 took $217$ and $199$ seconds at $N = 21$, W60 and W102 took
$6982$ and $6608$ --- and \fn{tests/test\_extend\_unitary.py} recomputes
\emph{every} stored unit of \emph{every} one of the sixteen rules from that
rule's own gate assignment, by label propagation on its own flip graph,
consulting no partner, and requires it to reproduce what both stores hold. Three
unrelated algorithms, rule by rule, $33$ units each. The identities below are
therefore results, and the figures draw both members of a pair from their own
data rather than mirroring one onto the other.

\paragraph{Reflection ($r_{01} \leftrightarrow r_{10}$) is a symmetry.} It leaves
the sector decomposition invariant for unitary rules at both boundary
conventions, so the sixteen rules give \textbf{twelve} distinct series and the
pairs $\{57,99\}$, $\{60,102\}$, $\{153,195\}$, $\{156,198\}$ agree term for
term. Table~\ref{tab:r19laws} is organised by these twelve classes.

\paragraph{The spin flip ($0 \leftrightarrow 1$ on every site) is not.}
$X^{\otimes N}$ is a basis permutation and $XHX$ has the same matrix magnitudes,
so it preserves the transition graph --- under \rt{pbc}. At \rt{obc0} the padding
is pinned to $|0\rangle$ and the flip sends it to $|1\rangle$. The consequence is
sharp: of the six spin-flip pairs, the only two whose series match are
$\{57,99\}$ and $\{156,198\}$, which are \emph{also} reflection pairs. The flip
never preserves a series on its own account (\fn{flip\_survives\_only\_via\_%
reflection}), and in two of the four remaining pairs it does something much
stronger than shift a curve:
\[
  \text{W60 (IIVV)}: N+1 \text{ sectors}
  \quad\longleftrightarrow\quad
  \text{W195 (VVII)}: 1 \text{ sector},
\]
and likewise W102 $\leftrightarrow$ W153. Same gate set, same lattice, same
brickwork; only the vacuum differs, and one of the pair fragments while the other
thermalises completely. The remaining two pairs stay within their class: W54
against W147 differ by exactly the one frozen vacuum state, and W108 against W201
is the familiar one-site shift.

\section{Every curve is an exact law}
\label{sec:laws}

\iftable{tab_r19_laws_obc0.tex}

None of the entries in Table~\ref{tab:r19laws} is a regression. Each was read
off the integer series and then evaluated at every $N$ in the sweep --- all
seventeen points, $N = 6 \ldots 22$, for every rule --- and the last column reports that
comparison, not a residual. \fn{scaling/paper\_figures.verify} performs it, the
notebook asserts on it before drawing, and \fn{tests/test\_paper\_figures.py}
runs it independently of either. Of the $16 \times 3 = 48$ (rule, series) slots,
$44$ carry an exact closed form; the four that do not are W156/W198's $D_{\max}$
and W108/W201's sector count, and the test suite asserts that the set of gaps is
exactly those four, so a curve cannot quietly lose its certificate.

Those four are covered by recurrences and an asymptote, checked the same way:
\begin{itemize}
  \item W108's sector count satisfies
    $a_N = 2a_{N-1} - a_{N-2} + a_{N-3}$ exactly, whose dominant root is
    $\rho^2 = 1.754878\ldots$, the square of the plastic number. R9 obtained
    that base by fitting and R18 derived it from the wall grammar; the
    recurrence is a third and purely arithmetic route to it.
  \item W201's sector count is W108's evaluated one site shorter,
    $n_\wcc^{201}(N) = n_\wcc^{108}(N-1)$, at every overlapping $N$.
  \item W156's $D_{\max}$ tag reads $4^{N/5}$ and is deliberately \emph{not} a
    law. At finite $N$ the largest sector is a mixture of rooms of length $2$ and
    $3$ (R2 sec.\,3), so no exact recurrence holds over the computed range. What
    is reported instead is the five-step ratio, which needs no fit:
    $\bigl(D_{\max}(21)/D_{\max}(16)\bigr)^{1/5} = 1.3195079$ against
    $4^{1/5} = 1.3195079$. Six digits by $N=21$ is evidence for the derived base;
    an exact-law assertion would be the wrong claim to make.
\end{itemize}

\subsection{$150/105$ is not one entry, and $60/102$ is}

The two polynomial tracks are usually written as pairs, but they are not pairs
in the same sense.

W60 and W102 agree \emph{term for term} in both series, so the figure labels
that track with an equals sign.

W105 and W150 do not. At $N\equiv 1 \pmod 4$ --- and, over the whole computed
range, at no other $N$ --- W105's largest sector carries \textbf{odd}
domain-wall number where W150's carries even ($N=9$: $\binom{10}{5} = 252$
against $\binom{10}{4} = 210$), and at exactly those $N$ W105 has one sector
fewer. It is one phenomenon appearing in both panels: an extra accessible level
of the wall number is precisely one fewer sector. The figure therefore labels
that track with a comma and draws it as two curves --- the second rule as a
larger open marker underneath the first, so that where the two coincide the
reader sees a filled mark inside a ring and where they part sees a lone ring.
The four lone rings in the sector panel are at $N = 9, 13, 17, 21$. Section
\ref{sec:frozen} adds a third face of the same parity obstruction.

\section{Figures 1 and 2 --- the scaling}

\begin{figure}[htbp]
\centering
\ifgraphic{../../figures/paper/fig_paper_scaling_obc0.pdf}{\linewidth}
\caption{Sector count (a) and largest sector (b) against chain length for all
sixteen unitary rules ($\rt{obc0}$, log $y$). Four rules fragment exponentially
--- W156${=}$W198 into $F_{N+2}$ sectors, W108 and W201 into $\Theta(\rho^{2N})$
--- while W60/W102 and W150/W105 produce only $\Theta(N)$ sectors and
correspondingly huge ones: W60's largest is exactly $2^{N-1}$, half the entire
Hilbert space. The grey and black tracks are the trivial extremes that bound
both panels: six ergodic rules with one sector of $2^N$ (W54: plus one frozen
vacuum state), and the identity W204 with $2^N$ sectors of one. Every tag is an
exact law (Table~\ref{tab:r19laws}) except W156's $4^{N/5}$, the derived
asymptote. Panels (a) and (b) are near mirror images because
$n_\wcc \cdot D_{\max} \ge 2^N$, and the trivial rules are exactly the ones for
which that holds with equality. All sixteen rules span the same grid
$N = 6\ldots22$.}
\label{fig:r19scaling}
\end{figure}

The two panels are also available separately as
\fn{fig\_paper\_sectors\_obc0.pdf} (Figure~\ref{fig:r19sectors}) and
\fn{fig\_paper\_dmax\_obc0.pdf} (Figure~\ref{fig:r19dmax}), at single-column
proportions, for a paper that wants them apart.

\begin{figure}[htbp]
\centering
\ifgraphic{../../figures/paper/fig_paper_sectors_obc0.pdf}{0.82\linewidth}
\caption{Number of Krylov sectors vs $N$ (\rt{obc0}), log scale --- the R2
Figure~1 content, rebuilt and completed. Exponential fragmentation ($156{=}198$,
$108$, $201$) against the linear and central-binomial families ($60/102$,
$150/105$), with the ergodic rules on the floor $n_\wcc = 1$ and the identity on
the ceiling $n_\wcc = 2^N$.}
\label{fig:r19sectors}
\end{figure}

\begin{figure}[htbp]
\centering
\ifgraphic{../../figures/paper/fig_paper_dmax_obc0.pdf}{0.82\linewidth}
\caption{Largest sector $D_{\max}$ vs $N$ (\rt{obc0}), log scale. The grey track
is the $2^N$ ceiling, traced exactly by the six ergodic rules (W54 one state
below it). The ordering is the reverse of Figure~\ref{fig:r19sectors}: the rules
with the fewest sectors have the biggest ones.}
\label{fig:r19dmax}
\end{figure}

\paragraph{Where each point comes from.} Not every rule had been pushed equally
far. The Tier-1a sweep reaches $N = 21$ on the fragmented rules and $22$ for
W150, but stops enlarging $N$ as soon as a rule is classified ergodic, so those
seven ended at $13$ there; the Tier-1e sweep covered all $256$ rules to $N = 16$
(a few to $17$). Since weak and strong components coincide for unitary rules the
two stores are merged, and \fn{wcc\_scc\_agreement} checks every overlapping unit
before the merge is used --- there is no disagreement anywhere in the sixteen.

That still left the census ragged, for a reason having nothing to do with the
rules, so \fn{qca\_fragmentation/extend\_unitary.py} closed it
(Section~\ref{sec:extend}). Each curve is nonetheless drawn only over its own
range and never extrapolated, and the machinery that annotates a short track's
$N_{\max}$ in the legend is still in place, because the ranges can diverge again
the moment anyone pushes one rule further.

\paragraph{On the drawing.} The $y$ axis is a true logarithmic axis with decade
ticks rather than $\ln(\cdot)$ on a linear axis, as in the original R2 figures.
A straight line still means an exponential, and the reader additionally gets to
read the actual sector counts off the axis. The five hues are the CVD-validated
categorical order used across the reports, re-checked over all pairs under
normal, protan, deutan and tritan vision; the worst separation is
$\Delta E = 7.4$ (W201 green against W60 blue, tritanopia), which is admissible
only with a secondary encoding, and three are present --- distinct markers, a
line style that tracks the class, and a direct law tag at every line end. The
green was darkened from \rt{\#008300} to \rt{\#006400} in this revision, which
raises that worst pair from $3.9$ and improves every other pair it takes part in.
The two trivial classes are drawn in neutrals rather than a sixth and seventh
hue, which is a statement about the content: they are the reference lines of the
figure, and the saturated hues belong to the rules that actually fragment.

\section{Closing the $N$ gap: \fn{extend\_unitary}}
\label{sec:extend}

The seven rules the sweeps stopped early on are the cheapest of the sixteen, not
the dearest. For a \emph{unitary} rule the sector partition is exactly the
connected-component structure of the single-flip graph
$x \sim x \oplus 2^i$ over the sites where the Hadamard fires --- a theorem, and
one validated against the streamed engine for all sixteen unitary rules at both
boundary conventions in \fn{tests/test\_flip\_graph.py}. Label propagation over
one \rt{uint32} array of $2^N$ entries decides it: $N = 22$ is $17$\,MB and a
quarter of a second per rule.

\fn{extend\_unitary} does that for the missing units and appends them to the
Tier-1e store in the ordinary schema, with \rt{checks.source} recording where
they came from. It writes nothing for a rule until the same code has reproduced
\emph{every} unit both stores already hold for it --- between $19$ and $28$ units
per rule, computed by union-find in one store and by Tarjan in the other, so
three unrelated algorithms have to agree before a fourth number is added. $93$
units were added this way, taking every unitary rule to $N = 22$. The whole run
takes well under a minute, most of it spent re-deriving the units that already
existed rather than on the new ones; run again it adds nothing and re-verifies
all $33$ units of each of the sixteen rules in half a minute.

Two things fell out of it. Every law in Table~\ref{tab:r19laws} continues to hold
at the new $N$, which is a real test rather than a restatement --- the extension
could have broken any of the $44$. And the fast path reproduces the Tier-1a
record for W150 at $N = 22$, which is the deepest unit in the whole Tier-1a
archive at either boundary condition --- every other rule stops at $21$ --- so
the check reaches the archive's own frontier.

The same command extends any unitary rule further:
\rt{python -m qca\_fragmentation.extend\_unitary -{}-n-max 24}. Memory is
$4 \cdot 2^N$ bytes, so $N = 26$ is $268$\,MB and $N = 30$ is $4.3$\,GB; the
reduction is proved only for unitary rules and refuses anything else.

\section{Figure 3 --- how big are the sectors?}

Neither scaling figure says anything about the \emph{distribution} of sector
sizes: $n_\wcc$ counts them and $D_{\max}$ reports one extreme. Nor does either
separate the two exponentially fragmented rules from each other --- both are
straight lines of similar slope. Figure~\ref{fig:r19hist} does both, and the two
columns turn out to look nothing alike.

\begin{figure}[htbp]
\centering
\ifgraphic{../../figures/paper/fig_paper_sizehist_obc0.pdf}{\linewidth}
\caption{Sector-size distribution of W156 (left) and W108 (right) at
$N = 13, 17, 21$ (\rt{obc0}), light to dark. \textbf{Top:} number of sectors per
log-spaced size bin. W156 is unimodal and peaked well away from $s=1$; W108 is
monotone decreasing over four decades. \textbf{Bottom:} the fraction of the
$2^N$ computational-basis states lying in sectors of size $\le s$ --- read the
median sector size off where a curve crosses the dotted line at $1/2$. The
dashed vertical rule is $D_{\max}$ at $N=21$; in both rules the typical sector
is far to the left of it. The sawtooth in the top-left panel is real and not
noise: at small $s$ a bin holds one or two integer sizes, and even sizes are
much commoner than odd.}
\label{fig:r19hist}
\end{figure}

\iftable{tab_r19_sizes_obc0.tex}

\paragraph{Two different kinds of exponential fragmentation.} W156 and W108 sit
close together in the growth-base plane --- both exponential in both
coordinates --- and Figure~\ref{fig:r19hist} shows that this hides a
qualitative difference. W156's sectors are all of comparable size: at $N=21$ its
$28\,657$ sectors have mean $73$ and median $108$, and only $12$ of them are
singletons. W108's are not: the same $N$ gives $170\,625$ sectors of which
$33\,552$, nearly a fifth, are frozen single states, the distribution falls
monotonically, and the median is $42$ against a $D_{\max}$ of $10\,946$.

\paragraph{$D_{\max}$ overstates the typical sector, increasingly.} The ratio
$D_{\max}/\text{median}$ runs $2.7 \to 3.2 \to 4.0$ for W156 over
$N = 13, 17, 21$ and $29 \to 76 \to 261$ for W108. Both grow, so the gap is not
a small-$N$ artefact. This matters wherever $D_{\max}$ stands in for ``how much
of the Hilbert space is dynamically connected'': for W108 at $N = 21$ the
largest sector holds $0.5\%$ of the space, but a state drawn uniformly at random
sits in a sector of about $42$ states out of two million.

\paragraph{Why this can be computed at all.} The stored \fn{sizes\_recurrent}
list is truncated at $2048$ entries for these rules once $N \ge 16$ (W108:
$N \ge 14$), which would cap Figure~\ref{fig:r19hist} at $N = 15$, or $13$ for
W108. The
\fn{size\_hist} field is \emph{not} truncated --- it carries every class --- and
\fn{size\_stats} asserts both that its class count equals $n_\wcc$ and that
$\sum_s s\,h(s) = 2^N$ before returning, so relying on it is checked rather than
assumed. That is what puts $N = 21$ within reach.

\section{Figure 4 --- the frozen states}
\label{sec:frozen}

Certifying the census produced a third exact series for free. The number of
size-one sectors is a closed form for \emph{every} unitary rule, and it is the
sharpest single number for telling the four classes apart.

\begin{figure}[htbp]
\centering
\ifgraphic{../../figures/paper/fig_paper_frozen_obc0.pdf}{0.82\linewidth}
\caption{Number of frozen states $n_1$ vs $N$ (\rt{obc0}), symlog $y$ so that
zero is a value rather than a gap --- which matters twice here: the six ergodic
rules have none at all, and W105 has none at $N \equiv 1 \pmod 4$. Exponential
rules have exponentially many (W108: a third of a million by $N=21$), the
polynomial rules have $O(1)$, and the identity has all $2^N$.}
\label{fig:r19frozen}
\end{figure}

Two of the laws are worth pulling out of Table~\ref{tab:r19laws}.

\paragraph{W108's frozen count factorises over the sublattices.}
\[
  n_1^{108}(N) \;=\; F_{\lfloor N/2\rfloor+2}\, F_{\lceil N/2\rceil+2},
\]
exactly, at every computed $N$ --- $25, 40, 64, 104, 169, \ldots$ up to
$54\,289 = F_{13}^2$ at $N = 22$. A state is frozen iff nothing can move on either
sublattice, and the two decouple, so the count is a product of two Fibonacci
numbers of nearly equal index. That is the same even/odd split that carries R18's
two independent domain-wall charges $|W_{\mathrm{even}}|$ and
$|W_{\mathrm{odd}}|$, seen here in a completely different observable. W201 obeys
the same form with the indices pulled down by one: $F_{m+1}^2$ at $N=2m$ and
$F_m F_{m+3}$ at $N = 2m+1$, the odd case being W108's value at $N-2$ off by
$(-1)^m$ via Catalan's identity.

\paragraph{W105 loses its frozen states exactly where it parts from W150.}
$n_1^{105}$ runs $1, 2, 1, 0$ as $N$ runs through $2, 3, 0, 1 \pmod 4$, so the
zeros fall at $N \equiv 1 \pmod 4$ --- precisely the residue class at which W105
has one sector fewer than W150 and its largest sector jumps to odd domain-wall
number. W150 never drops below one. Three observables, one parity obstruction.

\section{Using and editing the notebook}

Open \fn{notebooks/paper\_figures.ipynb} and run all cells; it writes into
\fn{figures/paper/} and takes a few seconds. Section~2 of the notebook holds
every adjustable parameter --- page widths, font sizes, the palette, the
\rt{TRACKS} list that decides which rules appear and how they are grouped, the
law tags and their nudges, the $y$ scale per series, the legend placement, the
histogram's $N$ values and bin count, the output directory. The drawing code
reads that cell and nothing else.

To add or regroup rules in the scaling figures, edit \rt{TRACKS}; rules inside
one track share a hue and the leading one is drawn filled on top of the others,
so a track that is secretly two curves gives itself away. To change which $N$ the
histograms show, edit \rt{HIST\_NS}. To change a \emph{law}, edit
\fn{scaling/paper\_figures.py} instead --- and expect the test suite to have an
opinion. The notebook itself is generated by
\fn{notebooks/\_build\_paper\_figures\_notebook.py}, so a hand edit that should
survive belongs in the builder.

\section{Caveats}

\begin{itemize}
  \item \textbf{\rt{obc0} only.} The \rt{pbc} sweep is held pending caveats
    (\fn{NEXT\_STEPS.md}); under \rt{pbc} at odd $N$ the brick-wall layers stop
    commuting, so those points would need even $N$ only. This matters more than
    usual for Section~\ref{sec:symmetry}: the ergodic-versus-fragmented contrasts
    it reports are \emph{created} by the \rt{obc0} vacuum, and under \rt{pbc} the
    spin flip is restored and the pairs would collapse onto each other.
  \item \textbf{The tail of every rule comes from the flip reduction}, not from
    the sweep that produced its head: $N \ge 17$ (or $18$) was computed by
    \fn{extend\_unitary}, and those records say so in \rt{checks.source}. The
    reduction is a theorem for unitary rules and for no others, it is checked
    against the engine independently, and the extension reproduces every stored
    unit of a rule before adding to it --- but it is a different algorithm from
    the one that produced the older points, and that is worth knowing.
  \item \textbf{These are WCC sectors}, which for the unitary rules shown
    coincide with the strongly connected components, so the decomposition is
    exact for the monitored and the unmonitored chain alike. That coincidence is
    a property of unitarity and does not survive dissipation --- see R-T14 on
    the recurrent-class/WCC distinction. It is also what licenses merging the
    Tier-1a and Tier-1e sweeps.
  \item \textbf{Both members of a reflection pair are measured and drawn}, as a
    filled marker inside a larger open one, rather than one being computed and
    the other inferred; a pair that ever stopped coinciding would show as a lone
    ring. Nothing plotted anywhere in this report is a symmetry image of another
    number.
  \item \textbf{Nothing here is new dynamics.} The sweep data is R2's and
    Tier-1e's, extended along the same $N$ axis; what is new is the completion of
    the census, the certification of the laws including the frozen-state series,
    the spin-flip statement, and the size distributions.
\end{itemize}

\end{document}
